\subsection{Native Coroutine Bridge}

Generator-based coroutines show the basic idea of cooperative suspension: execution runs until a suspension point, then control returns to an external caller. Native coroutines use a newer syntax and a more structured execution model, but the core idea remains related.

The main syntax is:

\begin{verbatim}
async def func(...):
    ...

await something
\end{verbatim}

An \texttt{async def} function does not behave like a normal function. Calling it creates a coroutine object. The body does not run immediately. The coroutine must be driven by an event loop or by another coroutine that awaits it.

\subsubsection{\texttt{async def} Functions as Native Coroutine Object Factories}

A normal function executes when called.

\begin{verbatim}
def normal_quote():
    print("normal body starts")
    return "D'oh!"

result = normal_quote()

print(f"result:       {result!r}")
print(f"type(result): {type(result)}")
\end{verbatim}

Possible output:

\begin{verbatim}
normal body starts
result:       "D'oh!"
type(result): <class 'str'>
\end{verbatim}

An \texttt{async def} function creates a coroutine object when called.

\begin{verbatim}
async def async_quote():
    print("async body starts")
    return "D'oh!"

coro = async_quote()

print(f"coro:       {coro}")
print(f"type(coro): {type(coro)}")

coro.close()
\end{verbatim}

Possible output:

\begin{verbatim}
coro:       <coroutine object async_quote at 0x...>
type(coro): <class 'coroutine'>
\end{verbatim}

The line \texttt{"async body starts"} was not printed. Calling \texttt{async\_quote()} created the coroutine object, but did not start executing the coroutine body.

The \texttt{close()} call is included here because the coroutine was created only for inspection. In real code, a coroutine should normally be awaited or scheduled.

To actually run it in a simple script, use \texttt{asyncio.run()}.

\begin{verbatim}
import asyncio

async def async_quote():
    print("async body starts")
    return "D'oh!"

result = asyncio.run(async_quote())

print(f"result:       {result!r}")
print(f"type(result): {type(result)}")
\end{verbatim}

Possible output:

\begin{verbatim}
async body starts
result:       "D'oh!"
type(result): <class 'str'>
\end{verbatim}

The expression:

\begin{verbatim}
async_quote()
\end{verbatim}

creates a coroutine object.

The expression:

\begin{verbatim}
asyncio.run(async_quote())
\end{verbatim}

creates a coroutine object, gives it to an event loop, runs it to completion, and returns the coroutine's final result.

A coroutine object can be inspected.

\begin{verbatim}
async def async_quote():
    return "No soup for you!"

coro = async_quote()

print(f"type(coro): {type(coro)}")
print()

selected_names = [
    "__await__",
    "send",
    "throw",
    "close",
    "cr_frame",
    "cr_code",
    "cr_running",
    "cr_await",
]

for name in selected_names:
    print(f"{name:<14} {name in dir(coro)}")

coro.close()
\end{verbatim}

Possible output:

\begin{verbatim}
type(coro): <class 'coroutine'>

__await__      True
send           True
throw          True
close          True
cr_frame       True
cr_code        True
cr_running     True
cr_await       True
\end{verbatim}

Some important attributes are:

\begin{itemize}
    \item \texttt{cr\_frame}: the coroutine's frame while it exists;
    \item \texttt{cr\_code}: the coroutine's code object;
    \item \texttt{cr\_running}: whether it is currently executing;
    \item \texttt{cr\_await}: what the coroutine is currently awaiting, if anything;
    \item \texttt{\_\_await\_\_}: the method that makes the object awaitable.
\end{itemize}

The coroutine's code object looks familiar.

\begin{verbatim}
async def async_quote():
    text = "The answer is 42."
    return text

coro = async_quote()
code_obj = coro.cr_code

print(f"type(code_obj): {type(code_obj)}")
print(f"co_name:        {code_obj.co_name}")
print(f"co_varnames:    {code_obj.co_varnames}")
print(f"co_consts:      {code_obj.co_consts}")

coro.close()
\end{verbatim}

Possible output:

\begin{verbatim}
type(code_obj): <class 'code'>
co_name:        async_quote
co_varnames:    ('text',)
co_consts:      (None, 'The answer is 42.')
\end{verbatim}

The coroutine object also has a frame before it completes.

\begin{verbatim}
async def async_quote():
    text = "Ni!"
    return text

coro = async_quote()

print(f"coro.cr_frame:       {coro.cr_frame}")
print(f"type(coro.cr_frame): {type(coro.cr_frame)}")

coro.close()
\end{verbatim}

Possible output:

\begin{verbatim}
coro.cr_frame:       <frame at 0x..., file "...", line 1, code async_quote>
type(coro.cr_frame): <class 'frame'>
\end{verbatim}

After the coroutine completes, its frame is gone.

\begin{verbatim}
import asyncio

saved_coro = None

async def async_quote():
    global saved_coro

    saved_coro = asyncio.current_task().get_coro()
    return "D'oh!"

result = asyncio.run(async_quote())

print(f"result: {result}")
print(f"saved_coro.cr_frame: {saved_coro.cr_frame}")
\end{verbatim}

Possible output:

\begin{verbatim}
result: D'oh!
saved_coro.cr_frame: None
\end{verbatim}

The completed coroutine no longer has an active frame.

The basic distinction is:

\begin{verbatim}
def function():
    calling it runs the body immediately

async def coroutine_function():
    calling it creates a coroutine object
    some scheduler must later drive that object
\end{verbatim}

\subsubsection{\texttt{await} as Structured Suspension Over Awaitable Objects}

The \texttt{await} expression suspends the current coroutine until another awaitable object completes.

It can appear only inside an \texttt{async def} function.

\begin{verbatim}
import asyncio

async def wait_and_quote():
    print("before await")

    await asyncio.sleep(0)

    print("after await")
    return "Nobody expects the Spanish Inquisition!"

result = asyncio.run(wait_and_quote())

print(f"result: {result}")
\end{verbatim}

Possible output:

\begin{verbatim}
before await
after await
result: Nobody expects the Spanish Inquisition!
\end{verbatim}

The line:

\begin{verbatim}
await asyncio.sleep(0)
\end{verbatim}

means:

\begin{verbatim}
create an awaitable sleep operation
suspend this coroutine until that operation completes
let the event loop run other ready work while this coroutine is suspended
resume this coroutine after the awaited operation completes
\end{verbatim}

The object returned by \texttt{asyncio.sleep(0)} is a coroutine object.

\begin{verbatim}
import asyncio

sleep_coro = asyncio.sleep(0)

print(f"sleep_coro:       {sleep_coro}")
print(f"type(sleep_coro): {type(sleep_coro)}")
print(f"'__await__' in dir: {'__await__' in dir(sleep_coro)}")

sleep_coro.close()
\end{verbatim}

Possible output:

\begin{verbatim}
sleep_coro:       <coroutine object sleep at 0x...>
type(sleep_coro): <class 'coroutine'>
'__await__' in dir: True
\end{verbatim}

An object can be awaited if it is awaitable. One practical test is whether it can provide an await iterator through \texttt{\_\_await\_\_}.

A coroutine may await another coroutine.

\begin{verbatim}
import asyncio

async def get_name():
    print("get_name starts")
    await asyncio.sleep(0)
    print("get_name resumes")
    return "Homer"

async def build_line():
    print("build_line starts")

    name = await get_name()

    print("build_line resumes")
    return f"{name}: 42"

result = asyncio.run(build_line())

print(f"result: {result}")
\end{verbatim}

Possible output:

\begin{verbatim}
build_line starts
get_name starts
get_name resumes
build_line resumes
result: Homer: 42
\end{verbatim}

The outer coroutine \texttt{build\_line()} suspends while waiting for \texttt{get\_name()} to finish. When \texttt{get\_name()} returns, its return value becomes the value of the \texttt{await} expression.

So this line:

\begin{verbatim}
name = await get_name()
\end{verbatim}

means:

\begin{verbatim}
run get_name() as an awaitable operation
pause build_line() until get_name() completes
bind get_name()'s returned object to local name name
\end{verbatim}

This resembles generator delegation, but it belongs to native coroutine scheduling rather than plain generator iteration.

A coroutine can keep local state across \texttt{await} suspension points.

\begin{verbatim}
import asyncio

async def counted_quote():
    count = 40

    print(f"before first await: count={count}")
    await asyncio.sleep(0)

    count = count + 1
    print(f"before second await: count={count}")
    await asyncio.sleep(0)

    count = count + 1
    print(f"before return: count={count}")
    return count

result = asyncio.run(counted_quote())

print(f"result: {result}")
\end{verbatim}

Possible output:

\begin{verbatim}
before first await: count=40
before second await: count=41
before return: count=42
result: 42
\end{verbatim}

The local variable \texttt{count} survives across suspension points. The coroutine frame is paused, not restarted.

The frame can be inspected while a coroutine is suspended by looking at task state from another coroutine.

\begin{verbatim}
import asyncio

async def worker():
    text = "D'oh!"

    await asyncio.sleep(0.01)

    return text

async def main():
    task = asyncio.create_task(worker())

    await asyncio.sleep(0)

    coro = task.get_coro()

    print(f"coro:       {coro}")
    print(f"type(coro): {type(coro)}")
    print(f"cr_frame:   {coro.cr_frame}")
    print(f"cr_await:   {coro.cr_await}")

    result = await task

    print(f"result:     {result}")
    print(f"cr_frame after completion: {coro.cr_frame}")

asyncio.run(main())
\end{verbatim}

Possible output:

\begin{verbatim}
coro:       <coroutine object worker at 0x...>
type(coro): <class 'coroutine'>
cr_frame:   <frame at 0x..., file "...", line 6, code worker>
cr_await:   <coroutine object sleep at 0x...>
result:     D'oh!
cr_frame after completion: None
\end{verbatim}

The \texttt{cr\_await} attribute shows the awaitable object currently being awaited.

A coroutine that awaits nothing will run straight through once the event loop starts it.

\begin{verbatim}
import asyncio

async def immediate():
    print("immediate starts")
    return 42

result = asyncio.run(immediate())

print(f"result: {result}")
\end{verbatim}

Possible output:

\begin{verbatim}
immediate starts
result: 42
\end{verbatim}

The function is still an \texttt{async def} function. Calling it still creates a coroutine object. It simply has no internal suspension point once scheduled.

The compact rule is:

\begin{quote}
\texttt{await} is a structured suspension point. It pauses the current coroutine until the awaited object completes, then resumes the coroutine with the completed result.
\end{quote}

\subsubsection{Event Loops as External Schedulers for Coroutine Progress}

A coroutine object does not advance itself. It needs a scheduler. In the standard library, the usual scheduler is an \texttt{asyncio} event loop.

The simplest script-level entry point is \texttt{asyncio.run()}.

\begin{verbatim}
import asyncio

async def main():
    print("main starts")
    await asyncio.sleep(0)
    print("main resumes")
    return "finished with 42"

result = asyncio.run(main())

print(f"result: {result}")
\end{verbatim}

Possible output:

\begin{verbatim}
main starts
main resumes
result: finished with 42
\end{verbatim}

The event loop drives the coroutine until completion. When the coroutine reaches \texttt{await}, the loop can run other ready work.

Two tasks can interleave cooperatively.

\begin{verbatim}
import asyncio

async def task(name, count):
    for i in range(1, count + 1):
        print(f"{name}: step {i}")
        await asyncio.sleep(0)

    return f"{name}: done"

async def main():
    task_a = asyncio.create_task(task("A", 3))
    task_b = asyncio.create_task(task("B", 2))

    result_a = await task_a
    result_b = await task_b

    print(result_a)
    print(result_b)

asyncio.run(main())
\end{verbatim}

Possible output:

\begin{verbatim}
A: step 1
B: step 1
A: step 2
B: step 2
A: step 3
A: done
B: done
\end{verbatim}

The exact interleaving can depend on scheduling details, but the important point is stable: each task gives control back to the event loop at \texttt{await asyncio.sleep(0)}.

The \texttt{Task} object wraps a coroutine and lets the event loop manage its progress.

\begin{verbatim}
import asyncio

async def worker():
    await asyncio.sleep(0)
    return "D'oh!"

async def main():
    task = asyncio.create_task(worker())

    print(f"task:       {task}")
    print(f"type(task): {type(task)}")
    print()

    selected_names = [
        "done",
        "result",
        "cancel",
        "get_coro",
        "add_done_callback",
        "__await__",
    ]

    for name in selected_names:
        print(f"{name:<20} {name in dir(task)}")

    result = await task

    print()
    print(f"result: {result}")
    print(f"task.done(): {task.done()}")

asyncio.run(main())
\end{verbatim}

Possible output:

\begin{verbatim}
task:       <Task pending name='Task-2' coro=<worker() running at ...>>
type(task): <class '_asyncio.Task'>

done                 True
result               True
cancel               True
get_coro             True
add_done_callback    True
__await__            True

result: D'oh!
task.done(): True
\end{verbatim}

A task is awaitable. Awaiting a task suspends the current coroutine until that task completes.

The coroutine inside the task can be inspected.

\begin{verbatim}
import asyncio

async def worker():
    await asyncio.sleep(0)
    return "No soup for you!"

async def main():
    task = asyncio.create_task(worker())
    coro = task.get_coro()

    print(f"coro:       {coro}")
    print(f"type(coro): {type(coro)}")
    print(f"cr_code:    {coro.cr_code}")
    print(f"cr_frame:   {coro.cr_frame}")

    result = await task

    print(f"result: {result}")

asyncio.run(main())
\end{verbatim}

Possible output:

\begin{verbatim}
coro:       <coroutine object worker at 0x...>
type(coro): <class 'coroutine'>
cr_code:    <code object worker at 0x..., file "...", line 3>
cr_frame:   <frame at 0x..., file "...", line 3, code worker>
result: No soup for you!
\end{verbatim}

The event loop does not make Python code run in parallel on multiple CPU cores. It coordinates progress between suspended coroutines.

A coroutine that never awaits will block other coroutines until it finishes.

\begin{verbatim}
import asyncio

async def greedy():
    print("greedy starts")

    for i in range(3):
        print(f"greedy working {i}")

    print("greedy ends")

async def polite():
    print("polite starts")
    await asyncio.sleep(0)
    print("polite resumes")

async def main():
    task_a = asyncio.create_task(greedy())
    task_b = asyncio.create_task(polite())

    await task_a
    await task_b

asyncio.run(main())
\end{verbatim}

Possible output:

\begin{verbatim}
greedy starts
greedy working 0
greedy working 1
greedy working 2
greedy ends
polite starts
polite resumes
\end{verbatim}

The \texttt{greedy()} coroutine does not suspend internally. Once it starts running, the event loop cannot move to another coroutine until \texttt{greedy()} returns.

This coroutine cooperates.

\begin{verbatim}
import asyncio

async def cooperative(name):
    for i in range(3):
        print(f"{name}: step {i}")
        await asyncio.sleep(0)

async def main():
    task_a = asyncio.create_task(cooperative("A"))
    task_b = asyncio.create_task(cooperative("B"))

    await task_a
    await task_b

asyncio.run(main())
\end{verbatim}

Possible output:

\begin{verbatim}
A: step 0
B: step 0
A: step 1
B: step 1
A: step 2
B: step 2
\end{verbatim}

Each \texttt{await asyncio.sleep(0)} gives control back to the event loop. The event loop can then resume another ready task.

The scheduling model is therefore still cooperative:

\begin{verbatim}
coroutine runs
    until it reaches await
        event loop regains control
        another ready coroutine may run
    awaited operation completes
        event loop resumes waiting coroutine
coroutine returns
    task becomes complete
\end{verbatim}

The event loop itself can be accessed from inside a running coroutine.

\begin{verbatim}
import asyncio

async def main():
    loop = asyncio.get_running_loop()

    print(f"loop:       {loop}")
    print(f"type(loop): {type(loop)}")

asyncio.run(main())
\end{verbatim}

Possible output:

\begin{verbatim}
loop:       <_UnixSelectorEventLoop running=True closed=False debug=False>
type(loop): <class 'asyncio.unix_events._UnixSelectorEventLoop'>
\end{verbatim}

The exact event loop class can differ by operating system and Python configuration. The conceptual role is the same: it is the external scheduler that drives coroutine progress.

A simple final example shows the complete path.

\begin{verbatim}
import asyncio

async def get_score(name):
    print(f"start get_score({name})")

    await asyncio.sleep(0)

    print(f"resume get_score({name})")
    return 42

async def build_report():
    print("build_report starts")

    score = await get_score("Homer")

    print("build_report resumes")
    return f"Homer score={score}"

result = asyncio.run(build_report())

print(f"result: {result}")
\end{verbatim}

Possible output:

\begin{verbatim}
build_report starts
start get_score(Homer)
resume get_score(Homer)
build_report resumes
result: Homer score=42
\end{verbatim}

The execution path is:

\begin{verbatim}
asyncio.run starts event loop
    build_report coroutine starts
        awaits get_score coroutine
            get_score starts
            awaits sleep
                event loop regains control
            sleep completes
            get_score resumes
            get_score returns 42
        build_report receives 42
        build_report returns report string
asyncio.run returns final result to normal synchronous code
\end{verbatim}

The bridge from generator-based coroutines to native coroutines is:

\begin{center}
\begin{tabular}{|p{0.28\linewidth}|p{0.34\linewidth}|p{0.28\linewidth}|}
\hline
\textbf{Mechanism} & \textbf{Generator coroutine} & \textbf{Native coroutine} \\
\hline
Function marker & \texttt{def} with \texttt{yield} & \texttt{async def} \\
\hline
Created object & generator object & coroutine object \\
\hline
Suspension syntax & \texttt{yield} & \texttt{await} \\
\hline
Manual driving & \texttt{next()}, \texttt{send()} & event loop / task / await \\
\hline
External scheduler & custom caller loop & \texttt{asyncio} event loop \\
\hline
Completion signal & \texttt{StopIteration.value} & returned result through \texttt{await} or \texttt{asyncio.run()} \\
\hline
\end{tabular}
\end{center}

The final rule is:

\begin{quote}
An \texttt{async def} call creates a coroutine object. \texttt{await} is the structured suspension point inside native coroutine code. An event loop is the external scheduler that resumes coroutines when their awaited operations are ready.
\end{quote}